\section{Start a new application}
\label{ch:quickstart}
\label{sec:own-data-workflow}

To start a new application, the analyst must construct model-ready node and
edge--mode flows. The package checks those inputs and computes the derivative;
it does not infer economic flows from raw vehicle counts, passenger counts,
schedules, or tonnage. The steps below define the policy, prepare the two input
tables, validate the baseline, and produce the welfare results.

\subsection{Check whether the method fits the policy}

The package is designed for marginal improvements to existing, positively
used edge--modes. The following screen should be applied before constructing
the input files.
\begin{longtable}{L{0.56\textwidth}L{0.36\textwidth}}
\toprule
Application & Implication \\
\midrule
\endhead
Marginal gross-benefit ranking of existing links with positive baseline use &
Directly supported. \\
Several small improvements to existing links &
Use a sparse policy bundle. \\
New link, new mode, or discontinuous topology change &
Outside the interior local derivative. \\
Large change to an existing link &
Use the derivative as a diagnostic and compare it with a nonlinear
counterfactual. \\
Capacity expansion &
State how capacity changes the congestion map; a free-flow cost reduction is
not automatically equivalent. \\
Project net present value or benefit--cost ratio &
Combine the welfare result with timing, costs, and omitted effects. \\
Distributional incidence, transfers, or toll revenue &
Add the relevant welfare weights and fiscal treatment. \\
\bottomrule
\end{longtable}

``Local'' describes the size of the policy perturbation. It does not describe
the adjustment horizon. In the economic-geography model, mobile labor
generally gives the derivative a long-run comparative-static interpretation.
The urban model likewise permits residence and workplace adjustment.

\subsection{Choose the economic setting and policy}

Begin by deciding which spatial response the application requires.
\begin{longtable}{L{0.22\textwidth}L{0.34\textwidth}L{0.34\textwidth}}
\toprule
Choice & Economic geography & Urban commuting \\
\midrule
\endhead
Movement & Goods between locations & Commuters between residences and
workplaces \\
Node margins & Labor and income & Residents and employment \\
Network flow & Value flow, in the same price basis as income & Commuter flow,
normalized by total commuters \\
Spatial response & Wages and mobile labor & Residence and workplace choices \\
Typical modes & Road, rail, and water freight & Road, bus, rail or streetcar,
and ferry \\
\bottomrule
\end{longtable}

Public transit is represented by bus, rail, streetcar, or ferry rows in the
urban edge--mode file. It does not require a separate transport solver.

Next state the policy. Select a mode, a proportional reduction in primitive
generalized cost, and one of three units: a directed edge, a reciprocal
physical link, or both. A physical-link experiment requires exactly two
opposite policy-mode edges with the same physical-link identifier. Unpaired
transit services remain directed policies.

\begin{center}
\small
\begin{tabularx}{\textwidth}{L{0.28\textwidth}L{0.23\textwidth}X}
\toprule
Feature & Status & Use \\
\midrule
Economic-geography derivative & Stable reference & U.S. freight and public
examples \\
\code{ChoiceLogsum} modes & Stable reference & Active multimodal
specification \\
Edge-local congestion & Stable reference & Headline freight specification \\
Fixed-route decomposition & Verified, memory intensive & Moderate networks
after the memory preflight \\
Endpoint-terminal congestion & Verified extension & Requires an
application-specific elasticity \\
Multimodal urban derivative & Verified extension & Requires commuting inputs
that satisfy the urban accounting identities \\
Seattle workflow & Candidate application & Additional ridership, cost, and
calibration evidence required \\
\code{ComponentCES} & Legacy only & Historical reproduction \\
\bottomrule
\end{tabularx}
\end{center}

\subsection{Install the package and initialize a project}

The package API is version 0.1.0 and requires Julia 1.10 or later. Continuous
integration tests Julia 1.10 and the current stable release on Linux, macOS,
and Windows. Git is needed to clone and update the repository. Running the
package requires Julia; building this PDF additionally requires
\code{make}, \code{latexmk}, Poppler, and \code{qpdf}. Release-maintainer
checks also use Python, but the grid analysis and guide build do not.

Clone the package and instantiate its committed Julia environment:
\begin{lstlisting}[language=bash]
git clone https://github.com/sfuchs-de/TransportNetworkWelfare.jl.git
cd TransportNetworkWelfare.jl
julia --project=. -e 'using Pkg; Pkg.instantiate()'
\end{lstlisting}
Keep each empirical application outside the package checkout. From the package
directory, create either an economic-geography or urban-commuting project:
\begin{lstlisting}[language=bash]
julia --project=. bin/tnw.jl init ../my-freight-network \
  economic_geography
julia --project=. bin/tnw.jl init ../my-commuting-network \
  urban_commuting
\end{lstlisting}
Omitting the final argument retains the economic-geography default. The command
writes the model-specific \code{config.toml} and seed CSV files, a
\code{sources.toml} provenance ledger, a preprocessing entry point, and a
project README:
\begin{lstlisting}
my-network/
  config.toml
  sources.toml
  data/nodes.csv
  data/edge_modes.csv
  scripts/prepare_inputs.jl
  figures/
\end{lstlisting}
It refuses to overwrite a nonempty directory. The seed inputs run immediately,
providing an installation check before the analyst replaces any data. The
preprocessing script marks where a reproducible raw-to-model-ready conversion
belongs, but it does not guess how vehicle counts, passenger records, or
schedules should be converted into the flows required by the model.

The self-contained examples ordinarily finish in under a minute after Julia
precompilation. Runtime depends on hardware and network size. Before
\code{decompose}, \code{validate} reports the estimated memory for the
fixed-route incidence matrix; ordinary \code{analyze} omits that dense matrix.
For the grid used in this guide, that estimate is
\GridDecompositionIncidenceGiB{} GiB.

\subsection{Replace the template inputs}

For economic geography, \code{nodes.csv} contains
\begin{lstlisting}
node_id,labor,income,longitude,latitude
\end{lstlisting}
For urban commuting, replace \code{labor} and \code{income} with
\code{residents} and \code{employment}. Coordinates are optional.

Each row of \code{edge\_modes.csv} represents one active mode on one directed
edge:
\begin{lstlisting}
edge_id,physical_link_id,origin,destination,mode,flow,
origin_terminal_id,destination_terminal_id,
congestion_elasticity
\end{lstlisting}
The final three columns are needed only for the corresponding congestion
specification. A missing mode row means that the mode is unavailable. Listed
flows must be finite and positive. Parallel physical routes require explicit
intermediate nodes; they cannot share the same ordered endpoint pair under
different edge identifiers.

The baseline must satisfy the accounting identities of the selected model.
Economic-geography inputs require
\[
 \sum_{j,m}\Xi_{ij,m}=\sum_{j,m}\Xi_{ji,m}.
\]
Urban inputs instead require
\[
 l_i^F+\sum_{j,m}\Xi_{ij,m}
 =
 l_i^R+\sum_{j,m}\Xi_{ji,m}.
\]
Construct any balancing, geographic matching, price conversion, or modal
rescaling in a separate preprocessing step. The generic adapter does not
perform those transformations silently.

Normalize economic-geography labor and income shares to one. Normalize urban
residence and workplace margins separately to one. Divide economic-geography
edge value flows by the same world-income aggregate and urban commuter flows
by the same total commuter count. Edge traffic itself generally sums to more
than one because one shipment or commuter can traverse several edges. Modal
shares sum to one within a directed edge.

\subsection{Convert standard source data into model inputs}

Most applications begin with several data products rather than the two tables
read by the package. A node geography supplies the unit of analysis, an
origin--destination table supplies residence--workplace or
production--consumption margins, a network file supplies available links, and
traffic or ridership evidence supplies baseline use. These sources must refer
to a common year, geography, population, and flow unit before they can be
combined. The table below gives common starting points and the construction
each one supports.

\begin{longtable}{L{0.21\textwidth}L{0.30\textwidth}L{0.41\textwidth}}
\toprule
Object & Common source & Use in the project \\
\midrule
\endhead
Location boundaries and centroids &
National statistical-office boundary files, such as ONS MSOA boundaries
\citep{ONS2021MSOABoundaries}, or a locally maintained GIS layer &
Choose node identifiers, calculate centroids, and provide the polygons used
for maps. Preserve the source coordinate system until the final export. \\
Residence--workplace flows &
Census or administrative workplace origin--destination tables, such as
ODWP01EW for England and Wales \citep{ONS2021WorkplaceFlows} &
Construct urban residence and employment margins. These tables describe OD
commuting, not the directed edge flows that the package reads. \\
Road and rail geometry &
Agency GIS files or OpenStreetMap extracts, including the downloadable
Greater London files distributed by Geofabrik
\citep{GeofabrikGreaterLondon} &
Define directed links, reciprocal physical-link identifiers, intermediate
route nodes, and map geometry. Record the extract date and license. \\
Transit service &
GTFS, timetable, and station-location feeds, such as TfL's open data
\citep{TfLOpenData} &
Define stops, links, service availability, travel times, and transfer
opportunities. A schedule is not a passenger-flow observation. \\
Traffic or ridership &
Count stations, ticketing records, passenger OD matrices, or freight-value
records; for London, TfL network-demand files and ORR station and OD data are
useful sources \citep{TfLNetworkDemand,ORRStationUsage2024} &
Allocate OD demand across directed edge--modes and check the resulting traffic
against independent counts. Station entries and exits alone do not identify
segment-level flow. \\
Generalized cost and congestion &
Timetables, speeds, fares, transfer penalties, capacity data, and local
engineering or econometric estimates &
Construct baseline cost components and choose elasticities. Keep values
transferred from another setting in a sensitivity range rather than presenting
them as local estimates. \\
\bottomrule
\end{longtable}

The package automates the calculation after these sources have been converted
into one internally consistent baseline. It automates schema checks, route
reconstruction, physical-link pairing, the welfare derivative, the
decomposition, sensitivity paths, manifests, and standard network figures. It
does not automate source-specific geocoding, the valuation of physical freight
or passenger counts, OD assignment, or the reconciliation of incompatible
populations. Those decisions belong in \code{scripts/prepare\_inputs.jl} and
must be described in \code{sources.toml}. A typical source entry records a
stable identifier, URL or local logical path, vintage, license, retrieval date,
and SHA-256:
\begin{lstlisting}
[sources.workplace_od]
title = "Census 2021 workplace flows, ODWP01EW"
url = "https://www.nomisweb.co.uk/sources/census_2021_od"
vintage = "2021"
retrieved = "YYYY-MM-DD"
license = "Open Government Licence"
sha256 = "replace-with-downloaded-file-hash"
\end{lstlisting}

\subsection{Prepare geometry and standard figures}
\label{sec:own-data-maps}

Coordinates in \code{nodes.csv} are enough for a schematic map with straight
links. For geographic maps, retain a polygon basemap and a line layer in the
same projected coordinate system as the node coordinates. The generic plotter
accepts GeoJSON. Each line feature must contain
\code{physical\_link\_id}; polygon and line features in the basemap need no
package-specific attributes. Shapefiles and GeoPackages can be used directly
in QGIS, or converted once with GDAL:
\begin{lstlisting}[language=bash]
ogr2ogr -f GeoJSON -t_srs EPSG:27700 \
  data/link_geometry.geojson raw/network_links.gpkg
ogr2ogr -f GeoJSON -t_srs EPSG:27700 \
  data/basemap.geojson raw/msoa_boundaries.gpkg
\end{lstlisting}
The CRS in these commands is illustrative. The analyst must choose one
appropriate to the application and use it consistently across nodes, links,
and polygons.

After \code{analyze} or \code{decompose}, create a welfare map with:
\begin{lstlisting}[language=bash]
python3 "$PKG/plots/network_example.py" \
  "$APP/data/nodes.csv" "$APP/data/edge_modes.csv" \
  "$APP/output/welfare_physical.csv" \
  "$APP/figures/welfare-map.pdf" \
  --metric primitive_F \
  --link-geometry "$APP/data/link_geometry.geojson" \
  --basemap "$APP/data/basemap.geojson"
python3 "$PKG/plots/hulten_vs_welfare.py" \
  "$APP/output/welfare_physical.csv" \
  "$APP/figures/traditional-versus-extended.pdf"
\end{lstlisting}
The map command also accepts a PNG output and \code{--transparent}. Without
\code{--link-geometry}, it draws straight links between node coordinates. The
second command draws the Traditional-versus-Extended comparison on common
axes and reports Pearson and rank correlations.
For publication maps that require labels, insets, or several policy layers,
join \code{welfare\_physical.csv} to the line layer by
\code{physical\_link\_id} in QGIS or another GIS package. The numerical output
remains the source of the color scale; the GIS file supplies geometry rather
than welfare data.

\subsection{Edit the configuration}

In the generated \code{config.toml}, the
\code{spatial\_specification} field selects the economic-geography or urban
commuting model. The configuration then supplies the corresponding spatial
parameters and route curvature.

Set \code{modal\_specification} to \code{choice\_logsum} and supply the modal
elasticity \(\eta\). The congestion block selects edge, endpoint-terminal,
composite, or no congestion. The remaining fields define the policy mode,
policy unit, shock fraction, output directory, data units, and all
transformations applied to the inputs.
Do not retain template parameter values merely because they pass validation.
Each value needs an empirical source or a stated sensitivity range. Appendix
\ref{app:config} gives the complete configuration and schema reference.

\subsection{Validate, analyze, and decompose}

For an application outside the package checkout, use explicit paths:
\begin{lstlisting}[language=bash]
PKG=/path/to/TransportNetworkWelfare.jl
APP=/path/to/my-network
julia --project="$PKG" "$PKG/bin/tnw.jl" validate "$APP/config.toml"
julia --project="$PKG" "$PKG/bin/tnw.jl" analyze "$APP/config.toml"
julia --project="$PKG" "$PKG/bin/tnw.jl" decompose "$APP/config.toml"
\end{lstlisting}
\code{validate} checks schemas, identifiers, accounting, route contraction,
modal interiority, policy pairing, congestion keys, regularity, and condition
numbers. Run \code{analyze} for the full welfare derivative. Run
\code{decompose} only when the reported fixed-route memory requirement is
feasible.

\begin{figure}[H]
\centering
\small
\begin{tabular}{c}
\fcolorbox{guideblue!45}{guideblue!4}{%
\parbox{0.82\textwidth}{\centering CSV/TOML inputs
\(\rightarrow\) schema, units, accounting, and route reconstruction}}\\[0.05em]
\(\downarrow\)\\[0.05em]
\fcolorbox{guideblue!45}{guideblue!4}{%
\parbox{0.82\textwidth}{\centering bilateral flows, modal and congestion maps
\(\rightarrow\) model-specific spatial Jacobian}}\\[0.05em]
\(\downarrow\)\\[0.05em]
\fcolorbox{guideblue!45}{guideblue!4}{%
\parbox{0.82\textwidth}{\centering equilibrium response and welfare projection
\(\rightarrow\) directed and physical-link welfare derivatives}}\\[0.05em]
\(\downarrow\)\\[0.05em]
\fcolorbox{guideblue!45}{guideblue!4}{%
\parbox{0.82\textwidth}{\centering model decomposition, diagnostics,
and hash-bound run manifest}}
\end{tabular}
\caption{Data-to-result workflow. Application-specific preprocessing precedes
the generic package; each package stage checks its output before the next stage
uses it.}
\label{fig:guide-pipeline}
\end{figure}

\subsection{Read and preserve the results}

\code{analyze} writes \code{welfare\_directed.csv} and, when requested,
\code{welfare\_physical.csv}. \code{decompose} writes the corresponding
\code{decomposition\_*.csv} files. The run also writes
\code{diagnostics.json} and \code{run\_manifest.json}.

A positive reported elasticity is a welfare gain from a primitive-cost
reduction. For a small reduction \(s\), the approximate percentage welfare
gain is \(100s\) times the elasticity. Compare \code{hulten} with
\code{primitive\_F} to see whether equilibrium adjustment changes the
traffic-only ranking. Interpret a decomposition only after its identity
residuals and condition diagnostics pass.

For example, the highest-ranked grid road link is
\GridTopLink. Its output row reports:
\begin{center}
\small
\begin{tabular}{lrrr}
\toprule
\code{physical\_link\_id} & \code{hulten} & \code{realized\_F} &
\code{primitive\_F} \\
\midrule
\GridTopLinkCode & \GridTopTraditional & \GridTopRealized &
\GridTopPrimitive \\
\bottomrule
\end{tabular}
\end{center}
Thus the Traditional approach gives the traffic-based elasticity
\GridTopTraditional, the full realized-cost effect is \GridTopRealized, and
the Extended approach gives the primitive-cost effect \GridTopPrimitive. A
one-percent primitive improvement implies
\(\Delta\log W\simeq\GridTopPrimitiveLogGain\), or approximately
\GridTopPrimitivePercentGain{} percent. Section \ref{ch:code} defines the
remaining output columns.

Retain the manifest for each result vintage. It records the package commit,
configuration, source-ledger, input, and output hashes together with the
parameters, transformations, and diagnostics.

\subsection{Planning an Elizabeth Line application}

Suppose the policy question is which operating Elizabeth Line segment would
yield the largest gross welfare gain from a small reduction in generalized
travel cost. This is an urban-commuting application because residence and
workplace margins differ and commuters can use road and public transport. It
is also a marginal application. The package can rank improvements to
positive-flow segments of the existing line or evaluate a bundle that improves
the operating corridor. It cannot recover the effect of opening the central
section from a zero-flow baseline. That question changes the available network
and requires a finite counterfactual.

Begin with the urban scaffold:
\begin{lstlisting}[language=bash]
PKG=/path/to/TransportNetworkWelfare.jl
APP=/path/to/elizabeth-line
julia --project="$PKG" "$PKG/bin/tnw.jl" \
  init "$APP" urban_commuting
\end{lstlisting}
Use 2021 MSOAs or another stable London geography as nodes. ONS ODWP01EW
reports counts by residence and workplace area; their row and column sums
provide resident and employment margins \citep{ONS2021WorkplaceFlows}. Census
Day fell during the
March 2021 lockdown, however, so these flows are not a post-opening commuting
baseline and must be treated as an initialization or replaced with more
representative OD evidence \citep{ONS2021TravelWorkQuality}. The source ledger
should state that limitation directly.

TfL timetable and station feeds identify the Elizabeth Line topology, travel
times, and stations, while the Greater London OpenStreetMap extract can supply
roads and detailed line geometry \citep{TfLOpenData,GeofabrikGreaterLondon}.
TfL station footfall and ORR station and OD statistics provide post-opening
evidence on use \citep{TfLNetworkDemand,ORRStationUsage2024}. None of these
sources alone is a complete segment-level commuting matrix. The preprocessing
script must reconcile their vintages, assign OD commuters to road and transit
paths, and report how closely the resulting edge flows reproduce the observed
station totals.

The model-ready urban node file contains normalized residents and employment.
The edge-mode file contains separate road and rail rows, directed commuter
flows normalized by the same total, and terminal identifiers where transfer
congestion is modeled. Set \code{mode = "rail"} in the policy block. Use
\code{unit = "both"} only where the two directions form a valid reciprocal
physical link; otherwise retain directed policies. The freight value
\(\eta=1.099\) can provide a transparent starting point for sensitivity, but
it is not a London estimate. A preferred application would estimate or
calibrate modal substitution from London evidence.

Before running the model, edit the following fields in place, retain the other
generated input and output entries, and declare the sensitivity values used in
the exercise:
\begin{lstlisting}
[input]
mode_order = ["road", "rail"]

[model]
modal_specification = "choice_logsum"
eta = 1.099

[congestion]
specification = "none"

[policy]
mode = "rail"
unit = "both"
shock_fraction = 0.01

[sensitivity]
eta = [0.75, 0.90, 1.099, 1.25, 1.40]
\end{lstlisting}
The no-congestion specification provides a transparent first run when no
London-specific congestion elasticity has been established. An empirical
application should add a congestion channel only after documenting its
definition, units, and calibration. It should also replace the template's
urban spatial parameters rather than treating them as London estimates. When
changing the generated composite specification to \code{none}, delete the
generated \code{[congestion.edge]} and \code{[congestion.terminal]} subtables.

\newpage
Once \code{scripts/prepare\_inputs.jl} writes the two CSVs, the full run is:
\begin{lstlisting}[language=bash]
julia "$APP/scripts/prepare_inputs.jl"
julia --project="$PKG" "$PKG/bin/tnw.jl" \
  validate "$APP/config.toml"
julia --project="$PKG" "$PKG/bin/tnw.jl" \
  analyze "$APP/config.toml"
julia --project="$PKG" "$PKG/bin/tnw.jl" \
  decompose "$APP/config.toml"
julia --project="$PKG" "$PKG/bin/tnw.jl" \
  sensitivity "$APP/config.toml"
python3 "$PKG/plots/network_example.py" \
  "$APP/data/nodes.csv" "$APP/data/edge_modes.csv" \
  "$APP/output/welfare_physical.csv" \
  "$APP/figures/elizabeth-line-welfare.pdf" \
  --metric primitive_F \
  --link-geometry "$APP/data/link_geometry.geojson" \
  --basemap "$APP/data/basemap.geojson"
\end{lstlisting}
The final map ranks marginal improvements to the represented operating
segments. A complete appraisal would add construction or operating costs,
reliability, crowding, emissions, distributional incidence, and any benefits
from new trips outside the maintained commuting baseline.

\subsection{Acceptance criteria}

The data are ready when their flow unit matches the selected spatial model and
congestion elasticities, and when the accounting and route-reconstruction
checks pass. The policy is ready when its direction, mode, reciprocal pairing,
and shock size describe the intended experiment. The model is ready when each
active parameter has a documented source, the calculation remains on a regular
branch, and finite-difference or limiting-case tests support any new extension.

The results are ready for interpretation when the sensitivity analysis
distinguishes changes in levels from changes in rankings. At that point, archive
the manifest with the model-ready inputs so the reported result can be
reconstructed from the same data and specification.

To update the package later, record the commit used for the archived results,
fast-forward the package checkout, instantiate the environment, rerun
\code{Pkg.test()}, and then revalidate the external project. Part II explains
the economics and algorithms behind each of these checks.
